\section{Discussion}
\label{sec:discussion}

\subsection{Why Can Weather-Free Models Beat an Operational Forecast?}
The official forecast of ONS uses weather forecasts and the experience of the operator's forecasting staff, information that our models do not have. Its lower accuracy in three of the four subsystems is therefore not a statement that simple models are better forecasters in general. Our results point to a specific mechanism. The official forecast's error is concentrated in the hours of PV generation, takes the form of an under-forecast of the global load, and grows over time together with the share of behind-the-meter MMGD (Fig.~\ref{fig:ons_bias}). The global load is the sum of the load measured at the transmission level and the estimated MMGD output; a forecast built on a component that does not keep up with the fast-growing installed MMGD capacity, or on profiles learned when this capacity was smaller, would under-forecast the midday global load in exactly this way. The models in this study learn the target directly from the most recent verified data: the level normalisation and the lagged profiles of days $d-2$, $d-7$, and $d-14$ make them adapt to the changing shape of the profile within a few weeks, without any explicit representation of PV. The fact that ENS shows no bias trend over almost three years, in all subsystems, supports this explanation, as does the reduction of the official forecast's bias from July 2026 onwards. That ENS remained more accurate after this reduction indicates that the timely use of the most recent load observations is a second, independent source of advantage. In the South and Southeast/Center-West, where the official forecast is competitive and more accurate in the afternoon, the weather information it contains is valuable, as also shown by the gains of the simple combination of the two forecasts.

These findings are relevant beyond Brazil. Many systems are going through a similar transition, and operators often forecast the load seen at the transmission level, net of behind-the-meter PV, or reconstruct the gross load with an estimate of PV output. In both cases, the forecast must keep up with installed capacity that changes every month. Our results suggest three simple practices for operators and forecasters: monitoring the bias of the forecast by hour of the day, not only its overall accuracy; maintaining at least one data-driven benchmark fitted on recent data, which is cheap and acts as an early-warning signal; and combining it with the operational forecast, which in our case reduced the error in the two subsystems where the official forecast is strongest.

\subsection{Practical Recommendations for Model Design}
For practitioners building day-ahead load models under drift, the experiments lead to concrete recommendations. (i) Normalise the target and the load inputs by a recent level before fitting tree-based models; without this, tree ensembles cannot follow a growing load, and the advantage of weekly refits or long windows is lost. (ii) Use calibration windows of three years or more: shorter windows cost more accuracy than any change in the refitting frequency. (iii) Frequent refitting is not necessary when the level is normalised: for windows of two years or more, monthly or even quarterly refits are within 0.05~pp of daily refits, which reduces the computational cost, especially for neural and tree models. (iv) Strong linear benchmarks should always be reported: ARX, which takes seconds to fit, was more accurate than a tuned GBM and an MLP, and the gain of the ensemble over ARX was small (0.04~pp) although significant in two subsystems. (v) For prediction intervals, empirical quantiles of recent out-of-sample errors are a better starting point than quantile regression with flexible models, and adaptive conformal inference with the operational delay restores the nominal coverage at little cost.

\subsection{Limitations}
This study has limitations that also indicate directions for further work. First, the models do not use weather forecasts. Historical numerical weather predictions for Brazil at the required resolution are not openly available for the whole period, and reanalysis data would give the models information unavailable at forecast time. Adding operational weather forecasts would most likely reduce the errors in the afternoon and in the South, where the official forecast is competitive. The weather-free setting, however, makes the comparison conservative with respect to our models. Second, the exact data cut-off of the official forecast is not documented. We used an issuance time of 12:00 on day $d-1$, consistent with the preparation of the daily operation programming, and showed that the conclusions for the North and Northeast hold even without observations of day $d-1$. Third, the verified global load contains an estimate of MMGD output produced by ONS, so the target itself is partly modelled. This affects the official forecast and our models in the same way, but the comparison should be interpreted as accuracy with respect to the operator's own verified series. Fourth, regional and municipal holidays, which matter in particular for the smaller subsystems, are not modelled. Finally, the evaluation covers one system. The mechanism we identify is general, but replicating the comparison with the data of other operators that publish their official forecasts, such as those of Australia or Europe, would confirm whether the size of the effect is similar elsewhere.
